\documentclass[12pt]{article}
\usepackage[utf8]{inputenc}
\usepackage[T1]{fontenc}
\usepackage{hyperref}
\title{Tinputz Language Resources for K-12 Curriculum Development}
\author{Compiled by Claude Sonnet 5 (claude-sonnet-5) for Andr\'as Kornai}
\date{}
\begin{document}
\maketitle

\section{Introduction}
Tinputz (ISO 639-3 tpz) is an Austronesian language of the Tinputz Rural LLG in north Bougainville, Papua New Guinea. It belongs to the Saposa--Tinputz branch of Northwest Solomonic. Its relatives are Teop, Hahon and Saposa. It has also been called Wasio, Buka and Kurtatchi. Reference works cite about 3,900 speakers (a 1991 figure), while the survey brief gives about 7,600. Many younger people also speak Tok Pisin or English. The searches found very little material aimed at schools.

\section{Primary Educational Resources}
No primers, textbooks or school materials for Tinputz turned up in the searches.

\section{Linguistic and Reference Materials}
\begin{itemize}
\item Wikipedia, Tinputz language (overview and classification): \url{https://en.wikipedia.org/wiki/Tinputz_language}
\item Wikipedia, Tinputz Rural LLG (the speech area): \url{https://en.wikipedia.org/wiki/Tinputz_Rural_LLG}
\item Encyclopedia.com, Kurtatchi (ethnographic entry): \url{https://www.encyclopedia.com/humanities/encyclopedias-almanacs-transcripts-and-maps/kurtatchi}
\end{itemize}

\section{Religious and Community Resources}
None were found.

\section{Digital and Audio Resources}
The only audio-related lead is a catalogue page of 1963 field recordings from north-east Bougainville at PARADISEC. The search result lists these under the names Timputz and Wasoi, and I did not check that the recordings are Tinputz.
\begin{itemize}
\item PARADISEC series NSBTIM: \url{https://paradisec.org.au/fieldnotes/NSBTIM.htm}
\end{itemize}

\section{Curriculum Development Context}
Almost nothing is available for Tinputz in the searches. A curriculum would have to be built from new documentation work. Any dictionary, grammar or other described material would need to be found by further searching in libraries and archives.
\end{document}